\chapter{Pruebas y validación}
\label{chap:testing}

La validación de LIA se ha planteado a partir de los riesgos del sistema: el
acceso indebido a datos de otra cuenta, la pérdida de coherencia durante la
ejecución de workflows, la aceptación de documentos no válidos y la
persistencia de respuestas de inteligencia artificial (IA) incorrectas o con
una estructura inesperada. Por este motivo, la estrategia combina pruebas
aisladas y rápidas con recorridos completos sobre infraestructura real. La
cobertura de código se utiliza como una señal auxiliar, pero no sustituye la
comprobación explícita de estos riesgos.

Este capítulo distingue dos actividades con objetivos diferentes. Las pruebas
automatizadas verifican que el software integra correctamente sus componentes,
valida los esquemas y responde de forma predecible ante errores. Por otra
parte, un banco experimental estudia la calidad, la latencia y el coste de
modelos generativos reales. En las pruebas automatizadas se utiliza un
proveedor determinista, por lo que su éxito no demuestra la calidad semántica
de un modelo externo.

\section{Estrategia y niveles de prueba}
\label{sec:testing-strategy}

La \Cref{fig:testing-levels} resume la organización de las pruebas. En la base
se encuentran las comprobaciones unitarias del dominio y de los casos de uso,
que permiten recorrer un número elevado de variantes sin acceder a servicios
externos. Sobre ellas se sitúan las pruebas de componentes y composables de la
interfaz, las integraciones con adaptadores y esquemas, las pruebas de extremo
a extremo de la API y, finalmente, los recorridos completos de navegador.

\begin{figure}[H]
  \centering
  \resizebox{\textwidth}{!}{%
    \begin{tikzpicture}[
  level/.style={
    draw=TFGBlue,
    rounded corners=3pt,
    align=center,
    minimum height=0.9cm,
    font=\small
  },
  note/.style={font=\scriptsize, text=TFGMuted, align=left}
]
  \node[level, fill=TFGLightBlue, minimum width=12.2cm] (unit) at (0,0) {
    Unitarias de dominio y aplicación · rápidas y aisladas
  };
  \node[level, fill=TFGLightBlue, minimum width=10.7cm] (component) at (0,1.25) {
    Componentes web, composables y adaptadores
  };
  \node[level, fill=TFGLightBlue, minimum width=9.2cm] (integration) at (0,2.5) {
    Integración con PostgreSQL, MinIO y esquemas
  };
  \node[level, fill=TFGLightBlue, minimum width=7.7cm] (api) at (0,3.75) {
    API E2E · NestJS + PostgreSQL
  };
  \node[level, fill=TFGLightBlue, minimum width=5.6cm] (browser) at (0,5) {
    Navegador E2E · recorrido crítico
  };
  \node[level, fill=TFGAccent!12, draw=TFGAccent, minimum width=4.8cm] (evaluation) at (5.5,5) {
    Evaluación IA real\\calidad, coste y latencia
  };

  \draw[-{Stealth[length=2mm]}, TFGBlue] (unit.north) -- (component.south);
  \draw[-{Stealth[length=2mm]}, TFGBlue] (component.north) -- (integration.south);
  \draw[-{Stealth[length=2mm]}, TFGBlue] (integration.north) -- (api.south);
  \draw[-{Stealth[length=2mm]}, TFGBlue] (api.north) -- (browser.south);
  \draw[-{Stealth[length=2mm]}, TFGAccent] (browser.east) -- (evaluation.west);

  \node[note, anchor=west] at (6.35,3.75) {Menos casos\\Mayor coste y realismo};
  \node[note, anchor=west] at (6.35,0.1) {Más casos\\Menor coste y mayor aislamiento};
\end{tikzpicture}

  }
  \caption{Niveles de prueba y evaluación experimental de IA.}
  \label{fig:testing-levels}
\end{figure}

La evaluación con modelos reales aparece como una rama separada porque no
persigue demostrar que una ruta HTTP funciona, sino estimar el comportamiento
de una dependencia probabilística. Esta separación permite ejecutar la suite
funcional sin credenciales, consumo facturable ni dependencia de cuotas
externas, y evita convertir una respuesta concreta del modelo en una aserción
frágil.

Los niveles se han distribuido de la siguiente forma:

\begin{itemize}
  \item \textbf{Pruebas unitarias de dominio y aplicación.} Verifican
  invariantes, transiciones de estado, autorización, aislamiento por cuenta,
  publicación de eventos y tratamiento de errores.
  \item \textbf{Pruebas de componentes y composables.} Comprueban el estado de
  sesión, los filtros, las mutaciones HTTP, la invalidación de caché y los
  controles visibles del workflow con la red simulada en su frontera.
  \item \textbf{Pruebas de integración.} Validan repositorios, serialización,
  adaptadores de almacenamiento y esquemas de salida de IA.
  \item \textbf{Pruebas E2E de API.} Arrancan la aplicación NestJS y una base de
  datos PostgreSQL reales; los servicios externos se sustituyen por dobles
  controlados.
  \item \textbf{Pruebas E2E de navegador.} Ejecutan Chromium contra la SPA y la
  API reales, con PostgreSQL y MinIO, para recorrer las operaciones críticas
  desde la perspectiva del usuario.
\end{itemize}

\section{Entorno y reproducibilidad}
\label{sec:testing-environment}

La campaña automatizada se ejecutó con Node.js 26.1.0 y pnpm 11.1.2. Las
pruebas E2E utilizaron PostgreSQL 16 y, cuando el recorrido incluía adjuntos,
MinIO. Cada base E2E se crea con un nombre reservado para pruebas y se regenera
antes de la ejecución, de modo que no se reutilizan datos de desarrollo ni de
producción. Los PDF empleados como fixtures carecen de información personal o
empresarial.

\begin{table}[H]
  \centering
  \small
  \caption{Entorno de ejecución de las pruebas.}
  \label{tab:testing-environment}
  \begin{tabular}{@{}p{3.4cm}p{4.4cm}p{6cm}@{}}
    \toprule
    \textbf{Nivel} & \textbf{Infraestructura} & \textbf{Dependencias controladas} \\
    \midrule
    Unitarias de API & Proceso Node.js & Repositorios, reloj, eventos y puertos mediante dobles. \\
    Componentes web & Vitest y \texttt{happy-dom} & Frontera HTTP y configuración de Nuxt simuladas. \\
    E2E de API & NestJS y PostgreSQL 16 & Almacenamiento en memoria y proveedor de IA determinista. \\
    E2E de navegador & Chromium, Nuxt, NestJS, PostgreSQL 16 y MinIO & Proveedor de IA determinista. \\
    Evaluación de IA & Adaptador Google y PDF ficticios & Prompts, esquemas, parámetros, casos y rúbrica congelados. \\
    \bottomrule
  \end{tabular}
\end{table}

La \Cref{tab:testing-environment} distingue la infraestructura real de los
dobles controlados en cada nivel. La integración continua reproduce la
instalación bloqueada de dependencias y
ejecuta comprobación de tipos, lint, pruebas unitarias, cobertura, construcción
y E2E de API. El E2E de navegador se mantiene como validación local porque
requiere el servicio MinIO y un navegador. Esta diferencia se registra de
forma expresa: la integración continua no se presenta como evidencia de un
recorrido que no ejecuta.

\section{Pruebas unitarias y de componentes}
\label{sec:unit-component-tests}

\subsection{Dominio y casos de uso del backend}

Las pruebas del backend se concentran en las reglas que pueden producir estados
incoherentes o accesos indebidos. En autenticación se verifican la creación de
sesiones, la rotación del token de refresco, la revocación ante reutilización,
el cierre de sesión y el cambio de cuenta activa. La administración de cuenta
incluye una comprobación específica de la guarda que admite a
\texttt{ADMIN} y rechaza a \texttt{MEMBER}.

En pipelines y workflows se prueban el orden y las restricciones de los
estados, la asignación y sustitución de definiciones, la creación de acciones,
el avance de pasos, las decisiones y las operaciones de completar, omitir y
reintentar. También se comprueba que una acción bloqueada o ajena al paso
actual no pueda modificarse. Estas pruebas son especialmente relevantes
porque una transición incorrecta afectaría a todas las vistas posteriores de
la oportunidad.

La cualificación verifica por separado preguntas de control, campos
personalizados y resúmenes. Los servicios de instanciación reutilizan la
instancia que ya pertenece a una oportunidad y rechazan plantillas de otra
cuenta. Las generaciones automáticas comprueban los estados pendiente,
completado y fallido, la persistencia de evidencia y la recuperación mediante
reintento. Para los adjuntos se validan la subida, descarga y eliminación, el
mantenimiento del contenido binario y la ausencia de la clave interna de
MinIO en los DTO expuestos.

El puerto de IA se sustituye en estas pruebas por un fake que devuelve objetos
deterministas conformes al esquema solicitado. Además, el adaptador de Google
se prueba de forma aislada para comprobar la construcción de la petición, el
envío de los buffers PDF y la propagación de metadatos. Con ello se valida la
integración técnica sin confundirla con una evaluación de exactitud documental.

\subsection{Componentes y lógica de la interfaz}

En el frontend se prueban las piezas con mayor densidad de comportamiento. El
store de autenticación cubre la inicialización de sesión, el cierre incluso
cuando falla la petición y el cambio de cuenta. El middleware comprueba la
redirección de usuarios autenticados y no autenticados. Los composables de
oportunidades verifican la serialización de filtros, las mutaciones, la subida
multipart de PDF y la invalidación de las consultas relacionadas.

Los componentes del workflow distinguen acciones pendientes, en curso,
fallidas, completadas, omitidas y bloqueadas. Las pruebas comprueban que cada
estado expone únicamente las operaciones permitidas. La red se simula en el
límite HTTP; de esta forma, la prueba sigue observando el comportamiento del
componente sin depender de que la API esté levantada.

\section{Pruebas de integración y E2E de API}
\label{sec:api-integration-tests}

Las pruebas de persistencia verifican que los repositorios conservan tanto el
valor funcional como los metadatos de generación: estado, evidencia y error.
Los adaptadores de almacenamiento comprueban que los buffers no se alteran y
que las operaciones de subida, descarga, firma y borrado se delegan
correctamente. El servicio común de MinIO elimina durante el arranque una
política pública previa y utiliza el tiempo de caducidad configurado para las
URL prefirmadas.

La suite E2E de API ejecuta doce casos con la aplicación NestJS y PostgreSQL
reales. Los casos más representativos se muestran en la
\Cref{tab:api-e2e-cases}.

\begin{longtable}{@{}L{4cm}L{7.7cm}L{2.1cm}@{}}
  \caption{Casos representativos de las pruebas E2E de API.}
  \label{tab:api-e2e-cases} \\
  \toprule
  \textbf{Riesgo} & \textbf{Comprobación} & \textbf{Resultado} \\
  \midrule
  \endfirsthead
  \multicolumn{3}{c}{\tablename\ \thetable{} -- continuación} \\[3pt]
  \toprule
  \textbf{Riesgo} & \textbf{Comprobación} & \textbf{Resultado} \\
  \midrule
  \endhead
  \bottomrule
  \endlastfoot
  Sesión y rutas protegidas & Registro, inicio de sesión, credenciales inválidas y rechazo de peticiones sin autenticación. & Superado \\
  Escalada de privilegios & Invitación, cambio de rol y expulsión permitidos a \texttt{ADMIN}; rechazo de \texttt{MEMBER} y de la autoeliminación. & Superado \\
  Validación de entrada & Rechazo de registros, oportunidades y documentos con forma o tipo no admitidos. & Superado \\
  Ejecución del proceso & Creación de oportunidad, asignación de workflow y finalización de una acción. & Superado \\
  Contexto documental & Almacenamiento de PDF y entrega de su buffer al proveedor determinista durante la generación. & Superado \\
  Aislamiento multiempresa & Rechazo de lectura o modificación de oportunidades, adjuntos, workflows e instancias de otra cuenta. & Superado \\
\end{longtable}

La ejecución concluyó con 12 casos superados de 12. La combinación de casos
positivos y negativos permite comprobar no solo que las operaciones válidas
funcionan, sino que las fronteras de cuenta y rol se aplican antes de modificar
el estado.

\section{Pruebas E2E de navegador}
\label{sec:browser-e2e-tests}

Los recorridos de navegador se implementan con Playwright. Sus aserciones para
la web reintentan automáticamente las condiciones visibles mientras la página
se actualiza, lo que permite observar el resultado final sin introducir esperas
fijas innecesarias \cite{playwrightAssertions2026}.

Se ejecutaron cuatro recorridos sobre Chromium:

\begin{enumerate}
  \item \textbf{Cualificación completa.} Registra un usuario, configura los
  elementos necesarios, crea una oportunidad, asigna su workflow, introduce
  información manual, genera un resumen con el proveedor fake y adjunta un PDF.
  Este caso también comprueba la descarga segura del documento.
  \item \textbf{Autenticación.} Muestra el error de credenciales no válidas y,
  después, establece una sesión real con las credenciales correctas.
  \item \textbf{Validación de formularios.} Comprueba los mensajes localizados
  de campos obligatorios y que un diálogo reabierto no conserva un error
  obsoleto.
  \item \textbf{Carga frente a estado vacío.} Distingue el fallo al recuperar
  pipelines de una respuesta válida sin elementos, evitando presentar ambos
  casos con la misma interfaz.
\end{enumerate}

Los cuatro recorridos finalizaron correctamente. El primero utiliza
PostgreSQL y MinIO reales, pero mantiene el proveedor de IA determinista. Por
tanto, demuestra la continuidad del flujo entre navegador, web, API, base de
datos y almacenamiento, no la calidad del texto generado por un modelo real.

\section{Validaciones de seguridad y robustez}
\label{sec:security-robustness-tests}

El aislamiento por \texttt{accountId} se comprueba en servicios, consultas y
E2E de API. Una cuenta no puede consultar o modificar la oportunidad de otra,
ni utilizar sus adjuntos, workflows o instancias de cualificación. La matriz de
roles verifica además que pertenecer a la cuenta no concede automáticamente
permisos administrativos.

La protección de documentos se prueba en el recorrido de navegador mediante
tres observaciones: la URL directa del objeto devuelve un rechazo, una URL
prefirmada vigente permite descargarlo y esa misma URL vuelve a ser rechazada
después de su caducidad. También se verifica que un fichero que no es PDF se
rechaza antes de alcanzar el almacenamiento.

Frente a fallos del proveedor de IA, los casos de uso persisten un estado de
error revisable y permiten reintentar la acción. Los esquemas estructurados
impiden aceptar como resultado funcional un objeto con una forma distinta de
la esperada. Estas medidas reducen errores de integración, pero no eliminan el
riesgo de que un resultado formalmente válido contenga información no
sustentada; dicho riesgo requiere la evaluación semántica descrita en la
\Cref{sec:ai-evaluation} \cite{autio2024}.

\section{Resultados de las pruebas automatizadas}
\label{sec:automated-test-results}

Todas las suites se regeneraron sin utilizar la caché de Turborepo. La
\Cref{tab:automated-test-results} reúne los resultados obtenidos.

\begin{table}[H]
  \centering
  \small
  \caption{Resultados de la validación automatizada.}
  \label{tab:automated-test-results}
  \begin{tabular}{@{}p{4.2cm}p{5.8cm}p{3.8cm}@{}}
    \toprule
    \textbf{Comprobación} & \textbf{Alcance} & \textbf{Resultado} \\
    \midrule
    API unitarias y componentes & 37 suites & 142 pruebas superadas \\
    Web unitarias y componentes & 11 archivos & 35 pruebas superadas \\
    API E2E & NestJS y PostgreSQL & 12 de 12 superadas \\
    Navegador E2E & Nuxt, NestJS, PostgreSQL, MinIO y fake de IA & 4 de 4 superadas \\
    Calidad estática & Tipos, lint y construcción & Superada; 13 avisos web no bloqueantes \\
    Integración continua & Tipos, lint, pruebas, cobertura, construcción y API E2E & Superada \\
    \bottomrule
  \end{tabular}
\end{table}

La API establece umbrales globales mínimos del 60\,\% para líneas y
sentencias, 37\,\% para funciones y 33\,\% para ramas. La web fija umbrales
inferiores acordes con el alcance actual de sus pruebas instrumentadas: 8\,\%
para líneas y ramas, 7\,\% para sentencias y 3\,\% para funciones. La
\Cref{tab:coverage-results} presenta los valores regenerados.

\begin{table}[H]
  \centering
  \small
  \caption{Cobertura de código por aplicación.}
  \label{tab:coverage-results}
  \begin{tabular}{@{}lrrrr@{}}
    \toprule
    \textbf{Aplicación} & \textbf{Líneas} & \textbf{Sentencias} & \textbf{Funciones} & \textbf{Ramas} \\
    \midrule
    API & 63,16\,\% & 62,53\,\% & 40,45\,\% & 34,59\,\% \\
    Web & 9,17\,\% & 8,25\,\% & 4,86\,\% & 9,60\,\% \\
    \bottomrule
  \end{tabular}
\end{table}

La cifra del frontend debe interpretarse junto con el resto de evidencias. La
cobertura instrumentada se concentra en middleware, stores, validación,
composables y componentes con lógica, mientras que cuatro recorridos E2E
atraviesan páginas completas. Esto no convierte el 9,17\,\% en una cobertura
suficiente para cualquier evolución: existe una deuda de regresión unitaria y
visual sobre páginas y estados secundarios. Sí demuestra que los riesgos
seleccionados no dependen exclusivamente de una métrica agregada.

\section{Evaluación experimental de inteligencia artificial}
\label{sec:ai-evaluation}

\subsection{Diseño del experimento}

El banco experimental se congeló antes de ejecutar los modelos. Contiene cinco
expedientes completamente ficticios, formados por siete PDF y 16 páginas. Los
casos varían en número de documentos, extensión, contenido accesorio e idioma,
y sus hashes SHA-256 permiten comprobar que no se modificaron durante la
campaña. Para cada expediente se realizan cuatro operaciones: resumen,
pregunta de control, campo personalizado y decisión de workflow.

Se planificaron tres repeticiones de cada operación para
\texttt{gemini-3-flash-preview} y \texttt{gemini-3.1-flash-lite}. Esto produce
60 salidas por modelo y 120 en total. Se fijaron prompts, esquemas, nivel de
razonamiento bajo, límite máximo de salida y política de reintentos. El
ejecutor conserva los resultados válidos y reanuda únicamente los pendientes,
una propiedad necesaria ante límites de cuota.

La evaluación semántica utiliza cuatro dimensiones puntuadas entre 0 y 2:

\begin{itemize}
  \item \textbf{exactitud}, ausencia de contradicciones con los documentos;
  \item \textbf{cobertura}, inclusión de los hechos relevantes esperados;
  \item \textbf{evidencia}, justificación mediante información localizable;
  \item \textbf{utilidad}, capacidad de apoyar la decisión sin contenido
  superfluo.
\end{itemize}

La puntuación máxima es 8. Dos agentes de IA independientes revisan copias
ciegas que ocultan el modelo, la repetición, las métricas y la valoración
ajena; no se presentan como evaluadores humanos. Tras congelar ambas
revisiones, los desacuerdos se reconcilian conservando las puntuaciones
originales y adoptando la valoración sustentada más estricta. Se considera
aceptable un modelo con una media mínima de 6/8, al menos un 95\,\% de salidas
estructuralmente válidas y ninguna alucinación crítica. Una alucinación es
crítica cuando una afirmación no sustentada altera la elegibilidad, el plazo,
el importe, una obligación o una decisión de continuar.

Además de la calidad, se registran validez estructural, latencia, tokens de
entrada y salida y coste estimado. El coste se calcula con la tarifa estándar
de pago vigente en la fecha del experimento, aunque las ejecuciones se
realizaran dentro de una cuota gratuita. Google publicaba 0,50 USD de entrada
y 3,00 USD de salida por millón de tokens para Flash, y 0,25 USD y 1,50 USD
para Flash-Lite, respectivamente \cite{googleGeminiPricing2026}.

\subsection{Resultados parciales y límite de cuota}

La cuota disponible permitió completar las 60 salidas de Flash-Lite y 32 de
las 60 previstas para Flash. Las 92 salidas obtenidas presentan una estructura
válida y cuentan con dos revisiones reconciliadas. La
\Cref{tab:ai-interim-results} muestra los resultados observados; las métricas
de Flash son provisionales porque no proceden todavía de la muestra completa.

\begin{table}[H]
  \centering
  \scriptsize
  \caption{Resultados parciales de la evaluación de modelos generativos.}
  \label{tab:ai-interim-results}
  \resizebox{\textwidth}{!}{%
    \begin{tabular}{@{}lrrrrrrr@{}}
      \toprule
      \textbf{Modelo} & \textbf{Salidas} & \textbf{Calidad} & \textbf{Válidas} & \textbf{Mediana} & \textbf{P95} & \textbf{Coste medio} & \textbf{Aluc. críticas} \\
      \midrule
      Gemini 3 Flash Preview & 32/60 & 7,75/8 & 100\,\% & 2152 ms & 8716 ms & 0,001169 USD & 0 \\
      Gemini 3.1 Flash-Lite & 60/60 & 7,53/8 & 100\,\% & 2663 ms & 5687 ms & 0,000800 USD & 2 \\
      \bottomrule
    \end{tabular}%
  }
\end{table}

Flash-Lite no satisface el criterio congelado, a pesar de su media y validez
estructural. Dos resúmenes del expediente ficticio de autobuses afirmaron que
el candidato acreditaba tanto una rampa como un espacio reservado, cuando el
documento citado únicamente acreditaba la rampa. La atribución convierte un
dato ausente en cumplimiento de un requisito y se clasifica como alucinación
crítica.

Flash obtiene de forma provisional una calidad media de 7,75/8 y no presenta
alucinaciones críticas en las 32 salidas disponibles. Sin embargo, comparar
ese subconjunto con las 60 salidas de Flash-Lite podría introducir un sesgo por
la combinación de casos y repeticiones que aún falta. Por ello no se selecciona
un modelo ganador ni se representa gráficamente la comparación. Las 28 salidas
restantes quedan pendientes de una futura reposición de cuota; esta limitación
se conserva como resultado del experimento, no se sustituye por estimaciones.

\section{Incidencias y limitaciones de la validación}
\label{sec:testing-limitations}

La campaña permitió identificar varias incidencias del propio entorno de
pruebas:

\begin{itemize}
  \item la suite E2E de Jest necesita finalizar de forma forzada después de
  cerrar NestJS y MikroORM, debido a manejadores que permanecen abiertos; los
  casos terminan correctamente, pero la causa debe eliminarse en una revisión
  futura del entorno;
  \item una primera ejecución de integración continua devolvió un error 500
  durante el registro porque el entorno E2E no definía
  \texttt{REFRESH\_TOKEN\_SECRET}; al conservar la causa original del error se
  identificó el problema, se añadió un secreto exclusivo de prueba y la
  ejecución completa pasó posteriormente;
  \item el lint del frontend mantiene 13 avisos no bloqueantes en componentes
  reutilizables de interfaz; no impiden la construcción, pero constituyen deuda
  de mantenimiento;
  \item el límite observado de solicitudes por minuto obligó a espaciar las
  peticiones y la cuota diaria impidió completar Flash, incluso con un ejecutor
  reanudable;
  \item los cuatro E2E de navegador cubren el recorrido crítico y varios estados
  de error, pero no sustituyen una campaña completa de regresión visual,
  accesibilidad, carga o compatibilidad entre navegadores.
\end{itemize}

Por tanto, la validación demuestra el funcionamiento de los recorridos
prioritarios, el aislamiento multiempresa, la autorización administrativa, la
seguridad básica de los adjuntos y la recuperación de errores de IA. No
demuestra todavía una selección definitiva de modelo, un rendimiento bajo
carga ni la ausencia de defectos en todas las combinaciones visuales.

\section{Trazabilidad de las pruebas}
\label{sec:testing-traceability}

La \Cref{tab:testing-traceability} relaciona los principales riesgos con los
requisitos definidos en el \Cref{chap:requirements}. La columna de evidencia
identifica el nivel más significativo, aunque un requisito pueda estar
cubierto por más de una suite.

\begin{longtable}{@{}L{3.6cm}L{3.1cm}L{5.4cm}L{1.6cm}@{}}
  \caption{Trazabilidad entre riesgos, requisitos y pruebas.}
  \label{tab:testing-traceability} \\
  \toprule
  \textbf{Riesgo} & \textbf{Requisitos} & \textbf{Evidencia principal} & \textbf{Estado} \\
  \midrule
  \endfirsthead
  \multicolumn{4}{c}{\tablename\ \thetable{} -- continuación} \\[3pt]
  \toprule
  \textbf{Riesgo} & \textbf{Requisitos} & \textbf{Evidencia principal} & \textbf{Estado} \\
  \midrule
  \endhead
  \bottomrule
  \endlastfoot
  Sesión inválida y acceso cruzado & RF-001--002; RN-001--002; RN-011 & Servicio de sesión y E2E de autenticación y aislamiento. & Completo \\
  Escalada de privilegios & RF-003; RN-003 & Guarda unitaria y matriz E2E de roles y autoeliminación. & Completo \\
  Estados incoherentes de proceso & RF-004--005; RF-009; RN-004--005; RN-007 & Entidades, handlers, eventos y recorridos E2E. & Completo \\
  Duplicación de cualificación & RF-006; RF-010; RN-006 & Servicios de instanciación y casos de uso de preguntas, campos y resúmenes. & Completo \\
  Salida de IA con forma incorrecta & RF-011; RF-015; RN-008; RNF-002; RNF-004 & Esquemas, fake determinista y adaptador de Google. & Completo \\
  Documento no válido o inseguro & RF-013--014; RN-002; RN-009; RN-011--012 & Handlers, API E2E y navegador E2E con acceso directo, firmado y caducado. & Completo \\
  Solvencia no sustentada & RF-012; RF-016; RI-008 & Banco ficticio, repeticiones y revisiones ciegas. & Parcial \\
  Error y recuperación de IA & RF-017--018; RN-007 & Persistencia de estado fallido, presentación del error y reintento de la acción. & Completo \\
\end{longtable}

Los resultados mantienen RF-012 y RF-016 como parciales. El sistema puede
generar evaluaciones revisables con contexto documental y se han obtenido
resultados experimentales, pero la comparación de modelos no está completa y
las métricas operativas no se persisten ni se explotan dentro de LIA. Esta
delimitación se retomará en el análisis de limitaciones y en la evaluación final
de los objetivos.
